\documentclass[11pt]{article}
% The preamble every project's paper shares: % The preamble every project's paper shares: % The preamble every project's paper shares: \input{../../preamble} from
% projects/<project>/paper/. Boards come from chessboard and xskak, plots from
% pgfplots; every number a paper prints comes from its generated.tex.
\usepackage[LSBC4,T1]{fontenc}
\usepackage[utf8]{inputenc}
\usepackage{lmodern}
\usepackage[margin=1in]{geometry}
\usepackage{microtype}
\usepackage{amsmath}
\usepackage{booktabs}
\usepackage{tabularx}
\usepackage{array}
\usepackage{xcolor}
\usepackage[obeyspaces,spaces]{url}
\usepackage{xskak}
\usepackage{pgfplots}
\pgfplotsset{compat=1.18}
\usepackage{caption}
\usepackage{subcaption}
\usepackage{enumitem}
\usepackage[section]{placeins}
\usepackage[hidelinks]{hyperref}

\setchessboard{boardfontencoding=LSBC4, boardfontfamily=skaknew, blackfieldcolor=black!22,
  whitefieldcolor=white, setfontcolors, showmover=false, label=true, labelfontsize=6pt,
  boardfontsize=14pt}

% A move or line in UCI, and a position in FEN, both breakable at a slash.
\newcommand{\uci}[1]{\texttt{#1}}
\newcommand{\fen}[1]{{\footnotesize\path{#1}}}
\newcolumntype{L}{>{\raggedright\arraybackslash}X}

\setlength{\parskip}{0.4em}
\setlength{\parindent}{0pt}
 from
% projects/<project>/paper/. Boards come from chessboard and xskak, plots from
% pgfplots; every number a paper prints comes from its generated.tex.
\usepackage[LSBC4,T1]{fontenc}
\usepackage[utf8]{inputenc}
\usepackage{lmodern}
\usepackage[margin=1in]{geometry}
\usepackage{microtype}
\usepackage{amsmath}
\usepackage{booktabs}
\usepackage{tabularx}
\usepackage{array}
\usepackage{xcolor}
\usepackage[obeyspaces,spaces]{url}
\usepackage{xskak}
\usepackage{pgfplots}
\pgfplotsset{compat=1.18}
\usepackage{caption}
\usepackage{subcaption}
\usepackage{enumitem}
\usepackage[section]{placeins}
\usepackage[hidelinks]{hyperref}

\setchessboard{boardfontencoding=LSBC4, boardfontfamily=skaknew, blackfieldcolor=black!22,
  whitefieldcolor=white, setfontcolors, showmover=false, label=true, labelfontsize=6pt,
  boardfontsize=14pt}

% A move or line in UCI, and a position in FEN, both breakable at a slash.
\newcommand{\uci}[1]{\texttt{#1}}
\newcommand{\fen}[1]{{\footnotesize\path{#1}}}
\newcolumntype{L}{>{\raggedright\arraybackslash}X}

\setlength{\parskip}{0.4em}
\setlength{\parindent}{0pt}
 from
% projects/<project>/paper/. Boards come from chessboard and xskak, plots from
% pgfplots; every number a paper prints comes from its generated.tex.
\usepackage[LSBC4,T1]{fontenc}
\usepackage[utf8]{inputenc}
\usepackage{lmodern}
\usepackage[margin=1in]{geometry}
\usepackage{microtype}
\usepackage{amsmath}
\usepackage{booktabs}
\usepackage{tabularx}
\usepackage{array}
\usepackage{xcolor}
\usepackage[obeyspaces,spaces]{url}
\usepackage{xskak}
\usepackage{pgfplots}
\pgfplotsset{compat=1.18}
\usepackage{caption}
\usepackage{subcaption}
\usepackage{enumitem}
\usepackage[section]{placeins}
\usepackage[hidelinks]{hyperref}

\setchessboard{boardfontencoding=LSBC4, boardfontfamily=skaknew, blackfieldcolor=black!22,
  whitefieldcolor=white, setfontcolors, showmover=false, label=true, labelfontsize=6pt,
  boardfontsize=14pt}

% A move or line in UCI, and a position in FEN, both breakable at a slash.
\newcommand{\uci}[1]{\texttt{#1}}
\newcommand{\fen}[1]{{\footnotesize\path{#1}}}
\newcolumntype{L}{>{\raggedright\arraybackslash}X}

\setlength{\parskip}{0.4em}
\setlength{\parindent}{0pt}

% Written by projects/learner-solve-chance/paper/generate.mjs; do not edit by hand.

\newcommand{\AppCommit}{\texttt{3a37748}}
\newcommand{\Draws}{1{,}000}
\newcommand{\Seed}{7}
\newcommand{\TargetCore}{75}
\newcommand{\TargetNew}{85}
\newcommand{\ThemedWin}{a fifth}
\newcommand{\ThemedLoss}{0.7}
\newcommand{\ObviousWin}{a tenth}
\newcommand{\ObviousLoss}{0.4}
\newcommand{\UnknownFirst}{half}
\newcommand{\ObviousUnknownWin}{a twentieth}
\newcommand{\NotesPoints}{60}
\newcommand{\TotalAttempts}{467}
\newcommand{\TotalCorrect}{441}
\newcommand{\Cards}{429}
\newcommand{\Once}{365}
\newcommand{\More}{64}
\newcommand{\FromEvents}{24}
\newcommand{\Rebuilt}{341}
\newcommand{\LogEntries}{106}
\newcommand{\LogOnce}{94}
\newcommand{\LogMore}{12}
\newcommand{\OnceMisses}{1}
\newcommand{\MoreMisses}{26}
\newcommand{\MoreWithMiss}{25}
\newcommand{\AllWithMiss}{26}
\newcommand{\RAttempts}{354}
\newcommand{\RDays}{21}
\newcommand{\RFirstTries}{365}
\newcommand{\RFrom}{2026-07-21}
\newcommand{\RTo}{2026-10-05}
\newcommand{\RSolved}{99.7}
\newcommand{\RHead}{52.0}
\newcommand{\RView}{78.1}
\newcommand{\RGap}{47.7}
\newcommand{\RGapLo}{40.5}
\newcommand{\RGapHi}{56.8}
\newcommand{\HOneDays}{11}
\newcommand{\HOneAttempts}{124}
\newcommand{\HOneView}{0.269}
\newcommand{\HOneHead}{0.694}
\newcommand{\HOneDiff}{$-$0.426}
\newcommand{\HOneLo}{$-$0.569}
\newcommand{\HOneHi}{$-$0.266}
\newcommand{\HTwoDays}{10}
\newcommand{\HTwoAttempts}{230}
\newcommand{\HTwoView}{0.270}
\newcommand{\HTwoHead}{0.779}
\newcommand{\HTwoDiff}{$-$0.509}
\newcommand{\HTwoLo}{$-$0.729}
\newcommand{\HTwoHi}{$-$0.348}
\newcommand{\PAttempts}{100}
\newcommand{\PDays}{seven}
\newcommand{\PFrom}{2026-08-28}
\newcommand{\PTo}{2026-10-05}
\newcommand{\PMisses}{four}
\newcommand{\PSolved}{96.0}
\newcommand{\PHead}{58.0}
\newcommand{\PView}{82.4}
\newcommand{\PRun}{93.1}
\newcommand{\PGap}{38.0}
\newcommand{\PGapLo}{31.3}
\newcommand{\PGapHi}{42.7}
\newcommand{\PViewGap}{13.6}
\newcommand{\PViewGapLo}{8.9}
\newcommand{\PViewGapHi}{22.1}
\newcommand{\PDiff}{$-$0.292}
\newcommand{\PDiffLo}{$-$0.366}
\newcommand{\PDiffHi}{$-$0.161}
\newcommand{\PHeadLL}{0.568}
\newcommand{\PHeadBrier}{0.191}
\newcommand{\PViewLL}{0.276}
\newcommand{\PViewBrier}{0.069}
\newcommand{\PRunLL}{0.298}
\newcommand{\PRunBrier}{0.064}
\newcommand{\HeadAbove}{60}
\newcommand{\ViewAbove}{287}
\newcommand{\ItemMean}{1131}
\newcommand{\PShiftHead}{500 or more (385 to 500, the grid's edge)}
\newcommand{\PShiftView}{305 (160 to 500, the grid's edge)}
\newcommand{\ShiftEdge}{500}
\newcommand{\SelfN}{40}
\newcommand{\SelfPred}{85.3}
\newcommand{\SelfObs}{92.5}
\newcommand{\SelfBrier}{0.075}
\newcommand{\SelfBaseBrier}{0.070}
\newcommand{\SelfLL}{0.288}
\newcommand{\SelfBaseLL}{0.295}
\newcommand{\ScoreTable}{%
\begin{tabular}{lrrr}
\toprule
estimate & mean chance & Brier & log loss \\
\midrule
headline rating (serving reads it) & 58.0\% & 0.191 & 0.568 \\
view of the themed sets & 82.4\% & 0.069 & 0.276 \\
running rate & 93.1\% & 0.064 & 0.298 \\
\midrule
share solved & 96.0\% & & \\
\bottomrule
\end{tabular}
}
\newcommand{\ReliabilityTable}{%
\begin{tabular}{lrrrrrr}
\toprule
& \multicolumn{3}{c}{headline rating} & \multicolumn{3}{c}{view} \\
\cmidrule(lr){2-4}\cmidrule(lr){5-7}
solve chance & attempts & mean chance & solved & attempts & mean chance & solved \\
\midrule
0.0 to 0.5 & 22 & 45.2\% & 90.9\% & 3 & 37.4\% & 100.0\% \\
0.5 to 0.6 & 37 & 55.2\% & 100.0\% & 0 & -- & -- \\
0.6 to 0.7 & 29 & 63.7\% & 93.1\% & 2 & 65.5\% & 100.0\% \\
0.7 to 0.8 & 8 & 74.2\% & 100.0\% & 17 & 75.8\% & 100.0\% \\
0.8 to 0.9 & 4 & 81.9\% & 100.0\% & 67 & 85.0\% & 94.0\% \\
0.9 to 1.0 & 0 & -- & -- & 11 & 92.1\% & 100.0\% \\
\bottomrule
\end{tabular}
}
\newcommand{\BucketTable}{%
\begin{tabular}{lrrrrr}
\toprule
bucket & attempts & mean rating & solved & headline's chance & view's chance \\
\midrule
fork & 11 & 906 & 100.0\% & 60.0\% & 84.0\% \\
hanging piece & 16 & 1511 & 87.5\% & 59.4\% & 82.4\% \\
mate in one & 37 & 1436 & 94.6\% & 49.4\% & 82.9\% \\
mate in two & 11 & 693 & 100.0\% & 61.4\% & 85.8\% \\
mate in three & 1 & 1075 & 100.0\% & 47.6\% & 27.0\% \\
planning & 24 & 716 & 100.0\% & 68.5\% & 81.5\% \\
\bottomrule
\end{tabular}
}
\newcommand{\LastDay}{2026-10-07}
\newcommand{\SyncedOn}{2026-10-13}
\newcommand{\ProfileCommit}{\texttt{d8b7712}}
\newcommand{\UpMin}{18}
\newcommand{\UpMax}{287}
\newcommand{\WithView}{six}
\newcommand{\WithoutView}{five}
\newcommand{\TotUpMin}{86}
\newcommand{\TotUpMax}{309}
\newcommand{\TodayTable}{%
\begin{tabular}{lrrrrrrr}
\toprule
& \multicolumn{2}{c}{headline} & \multicolumn{2}{c}{themed view} & \multicolumn{2}{c}{stored totals} & \\
\cmidrule(lr){2-3}\cmidrule(lr){4-5}\cmidrule(lr){6-7}
bucket & rating & aims at & rating & aims at & rating & aims at & hardest in the bank \\
\midrule
planning & 918 $\pm$ 222 & 725 & 1123 $\pm$ 94 & 930 & -- & -- & 1878 \\
hanging piece & 1608 $\pm$ 245 & 1415 & 1797 $\pm$ 108 & 1604 & -- & -- & 1899 \\
defensive move & 1096 $\pm$ 246 & 903 & none & -- & 1182 $\pm$ 70 & 989 & 1898 \\
mate in one & 1587 $\pm$ 247 & 1394 & 1833 $\pm$ 82 & 1640 & -- & -- & 1896 \\
fork & 1006 $\pm$ 280 & 813 & 1282 $\pm$ 119 & 1089 & -- & -- & 1897 \\
mate in two & 796 $\pm$ 290 & 603 & 1083 $\pm$ 118 & 890 & -- & -- & 1897 \\
mate in three & 1066 $\pm$ 314 & 873 & 1084 $\pm$ 212 & 891 & -- & -- & 1899 \\
pawn endgame & 1223 $\pm$ 318 & 1030 & none & -- & 1418 $\pm$ 112 & 1225 & 1887 \\
discovered attack & 1056 $\pm$ 328 & 863 & none & -- & 1362 $\pm$ 112 & 1169 & 1899 \\
London System & 1056 $\pm$ 328 & 863 & none & -- & 1362 $\pm$ 112 & 1169 & 1494 \\
f2 and f7 & 1040 $\pm$ 337 & 847 & none & -- & 1349 $\pm$ 85 & 1156 & 1889 \\
\bottomrule
\end{tabular}
}
\newcommand{\MateN}{37}
\newcommand{\MateSolved}{35}
\newcommand{\CurveHead}{(1,1332) (2,1335) (3,1337) (4,1337) (5,1338) (6,1345) (7,1346) (8,1349) (9,1351) (10,1375) (11,1357) (12,1360) (13,1369) (14,1396) (15,1397) (16,1402) (17,1403) (18,1419) (19,1421) (20,1429) (21,1430) (22,1433) (23,1441) (24,1469) (25,1471) (26,1491) (27,1494) (28,1496) (29,1497) (30,1514) (31,1521) (32,1540) (33,1542) (34,1548) (35,1548) (36,1557) (37,1560)}
\newcommand{\CurveView}{(1,1481) (2,1539) (3,1581) (4,1590) (5,1608) (6,1623) (7,1635) (8,1650) (9,1662) (10,1677) (11,1681) (12,1688) (13,1701) (14,1719) (15,1728) (16,1678) (17,1688) (18,1702) (19,1709) (20,1720) (21,1728) (22,1737) (23,1745) (24,1758) (25,1765) (26,1778) (27,1788) (28,1795) (29,1801) (30,1814) (31,1820) (32,1832) (33,1837) (34,1845) (35,1817) (36,1822) (37,1827)}
\newcommand{\ItemsSolved}{(1,1305) (2,1350) (3,1273) (4,1343) (5,1279) (6,1278) (7,1336) (8,1331) (9,1352) (10,1279) (11,1282) (12,1417) (13,1479) (14,1398) (16,1398) (17,1478) (18,1383) (19,1439) (20,1445) (21,1448) (22,1461) (23,1449) (24,1460) (25,1554) (26,1565) (27,1523) (28,1499) (29,1563) (30,1503) (31,1567) (32,1545) (33,1526) (35,1503) (36,1531) (37,1560)}
\newcommand{\ItemsMissed}{(15,1460) (34,1560)}
\newcommand{\ExRating}{1479}
\newcommand{\ExBucket}{mate in one}
\newcommand{\ExHead}{35}
\newcommand{\ExView}{78}
\newcommand{\ExHeadRating}{1369}
\newcommand{\ExViewRating}{1701}
\newcommand{\Example}{% Written by projects/learner-solve-chance/paper/generate.mjs; do not edit by hand.
\noindent\begin{minipage}[t]{0.37\textwidth}\vspace{0pt}
\chessboard[setfen={8/p1p2pkp/1p3p2/1Pb5/5Pq1/2P2R2/P3R2P/3B3K b - - 0 30}, boardfontsize=12pt, pgfstyle=straightmove, arrow=to, linewidth=0.07em, color=green!50!black, markmoves={g4-g1}]
\end{minipage}\hfill
\begin{minipage}[t]{0.6\textwidth}\vspace{0pt}\small
\textbf{The puzzle:} mate in one, rated 1479 on Lichess; Black to move.\par\smallskip
\textbf{The solution:} 30\ldots{}\symqueen{}g1\# (green), \uci{g4g1}.\par\smallskip
\textbf{Before it:} the headline rated him 1369 and gave 35\%; the view rated him 1701 and gave 78\%. He solved it.\par\smallskip
\textbf{FEN:} \fen{8/p1p2pkp/1p3p2/1Pb5/5Pq1/2P2R2/P3R2P/3B3K b - - 0 30}
\end{minipage}%
}


\title{Does the learner know how often he solves a puzzle?\\[0.3em]
\large The coaching app's solve chances, replayed over the player's puzzle record}
\author{Chess Coach research lab}
\date{}

\begin{document}
\maketitle

\begin{abstract}
The coaching app sets the difficulty of each puzzle set from a Glicko-2 rating: it aims where its
learner says the player will solve \TargetCore\% of the time. Replayed over the profile as the app
holds it (\PAttempts{} puzzle attempts on \PDays{} days, logged as he played), the rating serving
reads gave him a mean solve chance of \PHead\% on puzzles he solved \PSolved\% of. It sits close
to the difficulty he was served, so the sets aim far below him. The learner's view of his themed
sets, which it already keeps, gave \PView\%, still low, with about half the log loss. The test
committed before the run used a record rebuilt from the app's memory cards, which turned out to
keep almost only his solves, because a missed puzzle comes back for another review; its verdicts
do not stand. A second round, fixed before any of its attempts exists, tests the view on the
attempts the app records next. The app's owner adopted the view before it, on the profile's look;
round 2 decides whether it stays.
\end{abstract}

\section{The question}

The app's learner keeps a Glicko-2 rating for each skill and turns it into the chance that the
player solves an item of a given difficulty: Glicko's expected score against the item's rating,
with the item's own deviation. Serving reads that chance backwards. A bucket's puzzle set aims at
the difficulty he should solve \TargetCore\% of the time (\TargetNew\% while the bucket is new),
and draws the unseen puzzles nearest to it. The profile's puzzle totals count \TotalCorrect{}
solved of \TotalAttempts.

The learner keeps two estimates of each skill. The \emph{headline rating} folds every channel and
counts a result in an announced set at the Lichess discount, so that a record earned in announced
sets does not overstate how he plays when nothing is announced: in a bucket's own set a solve moves
it \ThemedWin{} as far as one in an unannounced set and a miss \ThemedLoss{} as far, and for an
obvious theme such as mate in N \ObviousWin{} and \ObviousLoss. The legacy attempt log cannot say
whether an attempt was a first try, so each of its attempts counts at \UnknownFirst{} of that. Beside it, a
\emph{view} per channel folds only that channel's first exposures, each at full weight. Serving
reads the headline. The learner's own notes expect the discount to read a player served at
\TargetCore\% in themed sets about \NotesPoints{} points under his strength. When the learner gives a solve
chance, how often does he solve? And is the rating serving reads the right one?

\section{The record and the rule}

\paragraph{The record.} The rule was committed before any solve chance was computed. It named a
record rebuilt so that every first try has its time. The app keeps a memory card for every puzzle
he has been served, with its last review's time and its passes and fails, and the bank gives each
puzzle's bucket and Lichess rating. A card reviewed exactly once holds one first try with its time
and result; a card reviewed more than once keeps only its last review's time, so it was left out.
That leaves \Once{} of the \Cards{} cards: \FromEvents{} with the app's own attempt event and
\Rebuilt{} rebuilt as the event the app writes today, \RFirstTries{} first tries from \RFrom{} to
\RTo.

\paragraph{The instrument.} For every first try, in time order, the app's own learner at commit
\AppCommit{} is folded on the record before it, at the moment of the attempt, and gives two solve
chances for the puzzle: from the headline rating and from the view of the attempt's channel (for a
bucket's own set, the themed view). An attempt is scored when its skill already has a rated attempt
in that channel, so both exist. A solve scores $-\ln p$ and a miss $-\ln(1-p)$, as the learner's
own self-check scores them.

\paragraph{The candidate.} Read a set's solve chance, and so its target, from the view of the set's
own channel, falling back to the headline where the view has no attempt yet. It fits nothing.

\paragraph{The rule.} The days with scored attempts alternate between two halves. (1) The headline
under-predicts if the share solved minus its mean chance has a 95\% interval over \Draws{}
bootstrap draws of the days (seed \Seed) entirely above zero, over-predicts if entirely below, and
is calibrated within the noise otherwise. (2) The view replaces the headline only if, on each
half, its mean log loss is below the headline's with the interval of the difference entirely below
zero. Named beside the verdicts, deciding nothing: the same replay on the profile as the app holds
it, scores against a running rate (the share of the skill's earlier attempts he solved), a
reliability table, the buckets, a best-fitting shift, the learner's self-check, and where each
estimate would aim a set today.

\section{Round 1, and why its verdicts do not stand}

On the rebuilt record (\RAttempts{} scored attempts on \RDays{} days) he solved \RSolved\%, and the
headline's mean chance was \RHead\%: a gap of \RGap{} points (\RGapLo{} to \RGapHi), so by rule 1
the headline under-predicts. By rule 2 the view passes on both halves: its log loss minus the
headline's is \HOneDiff{} on the first (\HOneDays{} days, \HOneLo{} to \HOneHi) and \HTwoDiff{} on
the second (\HTwoDays{} days, \HTwoLo{} to \HTwoHi).

Neither verdict stands. A missed puzzle becomes a redemption card that comes back until he solves
it, while a solved one retires, so a card reviewed once is nearly always a solve
(Table~\ref{tab:cards}). Leaving out the cards reviewed more than once removed \MoreWithMiss{} of
the \AllWithMiss{} puzzles he ever missed, which is the outcome both rules measure. The rebuilt
record's \RSolved\% is that selection, and on a record with almost no misses a higher chance can
only score better, so rule 2 could hardly fail. The memory a scheduler keeps is shaped by the results
it reacts to; a record of first tries has to come from a log written as they were played.

\begin{table}[!htbp]
\centering\small
\begin{tabular}{lrr}
\toprule
memory cards & cards & misses \\
\midrule
reviewed once (round 1's record) & \Once & \OnceMisses \\
reviewed more than once (left out) & \More & \MoreMisses{} on \MoreWithMiss{} cards \\
\bottomrule
\end{tabular}
\caption{Where the misses sit on the memory cards. A miss returns as a redemption card, so the
cards reviewed once hold almost none.}
\label{tab:cards}
\end{table}

\section{The profile as the app holds it}

The profile also holds what the app logged as he played: the legacy attempt log (\LogEntries{}
entries, \LogOnce{} of them on cards reviewed once and \LogMore{} on the others, all agreeing with
the cards in time, bucket, rating and result), the puzzle totals from before the log, which the
learner folds as a seed, and the attempt events. Every attempt in it was logged when it was played,
so nothing in it selects on the result. Replayed the same way, it gives \PAttempts{} scored attempts
on \PDays{} days, from \PFrom{} to \PTo, with \PMisses{} misses: he solved \PSolved\%.

The headline's mean chance was \PHead\%, a gap of \PGap{} points (\PGapLo{} to \PGapHi); the
view's was \PView\%, a gap of \PViewGap{} (\PViewGapLo{} to \PViewGapHi). Both under-predict, the
headline by far more (Table~\ref{tab:scores}). The view's log loss is about half the headline's,
a difference of \PDiff{} (\PDiffLo{} to \PDiffHi), and the reliability table shows where: the headline put
most of his attempts below a 70\% chance, and he solved nearly all of them
(Table~\ref{tab:reliability}). The same holds in every bucket with more than one attempt
(Table~\ref{tab:buckets}). The view scores only about as well as the running rate, though: the share of the skill's earlier
attempts he solved.

\begin{table}[!htbp]
\centering\small
\ScoreTable
\caption{Each estimate's mean solve chance and scores on the profile's \PAttempts{} scored
attempts. Lower is better for both scores.}
\label{tab:scores}
\end{table}

\begin{table}[!htbp]
\centering\small
\ReliabilityTable
\caption{The scored attempts binned by each estimate's solve chance.}
\label{tab:reliability}
\end{table}

\begin{table}[!htbp]
\centering\small
\BucketTable
\caption{Each bucket's scored attempts, the mean Lichess rating of its puzzles, the share solved
and each estimate's mean chance.}
\label{tab:buckets}
\end{table}

Why the headline reads him so low shows in its own fold. It starts each skill at the median rating
of the puzzles he was served, and a solve in a mate set that the log cannot call a first try moves
it \ObviousUnknownWin{} of a full step. So it follows the puzzles up and never gets ahead of them:
on the profile it sat on average \HeadAbove{} points above the puzzles he was asked about (their
mean rating \ItemMean), where Glicko gives a little over half, while the view sat \ViewAbove{}
points above. Figure~\ref{fig:mate} shows it on mate in one, where he solved \MateSolved{} of
\MateN. The shift that best fits the headline's chances, taken off every puzzle's rating, is
\PShiftHead; the view's is \PShiftView.

\begin{figure}[!htbp]
\centering
\begin{tikzpicture}
\begin{axis}[width=0.85\textwidth, height=0.45\textwidth, xmin=0, xmax=38,
  xlabel={his mate-in-one attempts, in order}, ylabel={rating},
  grid=major, legend cell align=left, legend columns=2,
  legend style={font=\scriptsize, at={(0.5,-0.22)}, anchor=north, column sep=6pt}]
\addplot[black, thick] coordinates {\CurveHead};
\addlegendentry{headline rating (serving reads it)}
\addplot[blue!70!black, dashed, thick] coordinates {\CurveView};
\addlegendentry{view of the themed sets}
\addplot[only marks, mark=*, mark size=1.3pt, gray] coordinates {\ItemsSolved};
\addlegendentry{puzzle solved, at its rating}
\addplot[only marks, mark=x, mark size=3pt, red!75!black, thick] coordinates {\ItemsMissed};
\addlegendentry{puzzle missed}
\end{axis}
\end{tikzpicture}
\caption{The two estimates before each mate-in-one attempt on the profile, and the puzzle's
Lichess rating. The headline follows the puzzles; the view sits well above them.}
\label{fig:mate}
\end{figure}

The learner's own self-check, which folds each channel's earlier attempts and predicts the later
ones, says the same of the views: on its \SelfN{} held-out attempts it predicted \SelfPred\% and he
solved \SelfObs\%, and its Brier score, \SelfBrier, is no better than its constant's,
\SelfBaseBrier. Figure~\ref{fig:example} shows the solved puzzle the headline gave the lowest
chance.

\begin{figure}[!htbp]\Example\caption{The solved puzzle on the profile the headline gave the
lowest chance, as he faced it, after the opponent's move.}\label{fig:example}
\end{figure}

\section{What it means for the app}

The sets aimed far below him. Today the view aims a set \UpMin{} to \UpMax{} points higher than the
headline in the \WithView{} buckets that have a themed view (Table~\ref{tab:today}). The other
\WithoutView{} have no attempt in their own sets since the log began, and the app aims them from
their stored totals read alone at full weight, \TotUpMin{} to \TotUpMax{} points higher than the
headline. The bank can serve every one of those targets. The view still under-predicts him, so sets
aimed from it should still land above \TargetCore\%, but nearer.

By its rule this round adopts nothing: its test did not stand, and the profile's look, named
beside it, decides nothing. On that look the app's owner chose to adopt the view at once: since
chess PR 78 the app aims a bucket's set from its themed view, and the skill map's next difficulty
reads the same estimate; since PR 79, a bucket with no themed view aims from its stored totals
instead of the headline. Commit \AppCommit{}, which this paper measures, has both. Round 2 decides
whether the view stays; it scores no attempt in a bucket without one. It takes the attempts the app
records after the last one in the profile at commit \AppCommit{} (played on \LastDay, Los Angeles
time), each folded on everything before it as the app holds it, and runs once they fall on at least
ten days. Rules 1 and 2 apply unchanged: sets now aim harder, and each estimate is still scored on
the puzzles he was served. A third rule, fixed after round 2's first attempts were recorded and
before any of their results was read, tests whether the chance holds for one-move puzzles and longer
ones alike, which a target per depth would need; the project's README states it.

\begin{table}[!htbp]
\centering\footnotesize\setlength{\tabcolsep}{4pt}
\TodayTable
\caption{Each puzzle bucket as the profile stood at its last sync, on \SyncedOn{} (Los Angeles
time): each estimate's rating with its deviation, the difficulty a set
would aim at from it (the core target, clamped to what the bucket can serve), and the hardest
puzzle the bank holds. The stored totals, read alone, are shown where the app aims from them: in a
bucket with no themed view.}
\label{tab:today}
\end{table}

\section{Limitations}

\begin{itemize}[nosep]
\item One player, served puzzles far below him: the profile holds \PMisses{} misses, so its log
losses rest on a handful of them.
\item The profile's attempts fall on \PDays{} days, so the bootstrap over days has few to draw from
and its intervals are rough.
\item The log cannot say which attempts were retries; the replay counts them as the app does, at
half weight in the headline and in full in the view.
\item A puzzle's published Lichess rating stands for its difficulty, as in the app. An item-level
model would need more than one attempt per puzzle.
\item Round 1's record was rebuilt from memory that selects on the result. Round 2 reads only the
app's own log, written as the attempts are played.
\end{itemize}

\end{document}
